\documentclass[11pt]{article}

\usepackage{graphicx} 
\usepackage{geometry}
\geometry{
    left=1.5cm,
    right=1.5cm,
    top=2cm,
    bottom=2cm
}
\usepackage[english]{babel}

\usepackage{xcolor}
\definecolor{lightblue}{RGB}{182, 249, 255}
\definecolor{darkblue}{RGB}{19, 75, 95}
\definecolor{bluedot}{RGB}{81,171,203}

\usepackage{hyperref}
\hypersetup{
    colorlinks=true,
    linkcolor=blue,
    filecolor=magenta,      
    urlcolor=darkblue,
    citecolor=blue,
    pdftitle={CitationCheck documentation}
}
\urlstyle{same}

\begin{document}
    \begin{center}
        \includegraphics[scale=0.5]{../assets/logo.png}\vspace{5mm}

        {\Huge \textbf{CitationCheck}}\vspace{2mm}

        {\Large v1.0}\vspace{2mm}

        {\Large \texttt{Documentation}}
    \end{center}
    \hrule \vspace{10mm}

    \section{Installation}

        \subsection{Windows}

        On Windows, please find the executable and start it.
    
        \subsection{Other}

        I plan to support other operating systems as well. For now, simply use the \textbf{PY}-file to start it. Moreover, 
        due to the fact that I do not possess a MacBook, I could not test whether the program works correctly, although it should since Python is platform-independent.

    \section{Data}

    The data about retractions is collected from the service \textit{RetractionWatch}. It is retrieved using \textit{Crossref}.
    Since the application fetches data via the public \textit{Crossref}-API, the amount of requests that the API receives is unpredictable.
    Therefore, \textit{Crossref} recommends to include an email address in the requests. The requests will be directed to a special pool of servers that are more reliable.
    Your data is not used for spam and it will not be sold to anyone. For more information, click \href{https://www.crossref.org/documentation/retrieve-metadata/rest-api/tips-for-using-the-crossref-rest-api/}{here}.

    \section{Usage}

        \subsection{Start}
    
        Start the executable on your computer. If there is no executable, you can also download the repository (see section \ref{sec:help}) and execute the python script.

        \subsection{BiB or SINGLE}

        When the GUI shows up, you can choose between two options: \textbf{BiB} and \textbf{SINGLE}.

        \begin{itemize}
            \item \textbf{[BiB]}
            
            The \textbf{BiB}-option opens a window where you can choose your \textit{bibtex}-file. The process starts immediately after you have selected a file.
            
            \item \textbf{[SINGLE]}
            
            The \textbf{SINGLE}-option opens a window where you can enter a title and a DOI. It only performs a search for one reference.
            The process starts immediately after you have confirmed the input.
        \end{itemize}

        \subsection{CSV or API}

        The GUI lets you choose between another two options: \textbf{CSV} and \textbf{API}.

        \begin{itemize}
            \item \textbf{[CSV]}
            
            The \textbf{CSV}-option tells the application to use the provided CSV-file. Although the CSV file contains the same information that can be retrieved via the API,
            it is not held up-to-date. To ensure this, press the \textbf{UPDATE}-button to download the latest version. Depending on your internet connection, it will take some time (e.g. 10 seconds (as of now the file has a size of roughly 46 MB))
            
            \item \textbf{[API]}
            
            The \textbf{API}-option uses the \textit{Crossref}-API to retrieve data. An internet connection is required for this.
        \end{itemize}

        \subsection{Results}

        The two fields \textit{passed} and \textit{failed} show how many references have passed the check and how many did not. If you would like to have a more detailed report, 
        click on the \textbf{PDF}-button and save the pdf. 

    \section{Speed}
    
    The time it takes for the application to present a result varies greatly. When you use the API, it depends on the quality of your internet connection and how busy the \textit{Crossref}-API is.
    When you use the CSV file, it depends on the computing power of your computer.
    Plus, the average amount of references per paper varies depending on the field and the topic. Nevertheless, to give a you a number, a bibtex-file with $50$ citations roughly needs $47$ seconds (CSV-option).
    You can assume that $1$ reference takes about $1$ second.

    \section{Help}\label{sec:help}

    The project's homepage can be found on GitHub using this link: \\ \textbf{\url{https://github.com/BenSt099/CitationCheck}}. \vspace{5mm} 
    
    \noindent If you want to report a bug, a feature or run into an issue, you can create an \textit{issue} on GitHub: \\ \textbf{\url{https://github.com/BenSt099/CitationCheck/issues}}.

\end{document}